\subsection{半单环}

定理 1（韦德伯恩）. —— 设 A 是一个环。下列性质等价：

\begin{enumerate}
\def\labelenumi{(\roman{enumi})}
\item
  A-模 \(\mathrm{A}_s\) 是半单的。
\item
  对 A 的每个左理想 \(\mathfrak{a}\)，存在 A 的一个左理想 \(\mathfrak{b}\)，使得 \(\mathrm{A}_s\) 是 \(\mathfrak{a}\) 与 \(\mathfrak{b}\) 的直和。
\item
  环 A 是左阿廷的，且 (A,A)-双模 \(_{s}A_{d}\) 是半单的。
\item
  环 A 同构于一个有限族单环的乘积。
\item
  存在一个整数 \(s \geqslant 0\)、诸体 \(\mathrm{D}_1, \ldots, \mathrm{D}_s\) 以及整数 \(r_1 \geqslant 1, \ldots, r_s \geqslant 1\)，使得环 A 同构于诸矩阵环 \(\mathbf{M}_{r_i}(\mathrm{D}_i)\) 的乘积。
\item
  环 A 是左阿廷的，且存在一个忠实半单的左 A-模。
\end{enumerate}

(i) 与 (ii) 的等价性由 VIII, 第 56 页的推论 2 得出，而 (iv) 与 (v) 的等价性由 VIII, 第 120 页的定理 1 得出。

A-模 \(A_{s}\) 是有限生成的。若它是半单的，则它是左阿廷的（VIII, 第 71 页，命题 10）。由于 A-模 \(A_{s}\) 是忠实的，这就证明了 (i) 蕴含 (vi)。反之，设性质 (vi) 成立；设 M 是一个忠实半单的 A-模。存在一个整数 \(m \geqslant 1\)，使得 \(A_{s}\) 同构于 \(M^{m}\) 的一个子模（VIII, 第 50 页，命题 5, a)）；由于 M 是半单的，\(A_{s}\) 亦然。于是我们证明了 (i) 与 (vi) 的等价性。

证明 (i) 蕴含 (iii)。设环 A 具有性质 (i)；我们已注意到此时 A 是左阿廷的。现在，左 A-模 \(A_{s}\) 的自同态就是 A 中元素的右乘。(A, A)-双模 \(_{s}A_{d}\) 是半单的这一事实随即由 VIII, 第 86 页的命题 6 得出。

证明 (iii) 蕴含 (iv)。设 (A, A)-双模 \(_{s}A_{d}\) 是半单的。它是有限生成的，故存在单 (A, A)-子双模的一个有限族 \((\mathfrak{a}_{i})_{i\in I}\)，其直和为 \(_{s}A_{d}\)。换言之，诸 \(a_{i}\) 是 A 的非零双边理想，A 的加群是诸 \(a_{i}\) 的直和，且对每个 \(i\in I\)，A 的含于 \(a_{i}\) 的每个双边理想都等于 0 或 \(a_{i}\)。令 \(b_{i}=\sum_{j\neq i}a_{j}\)（对每个 \(i\in I\)）；这是 A 的一个双边理想。映射 \(a\mapsto(a+\mathfrak{b}_{i})_{i\in I}\) 是从环 A 到诸环 \(A/b_{i}\) 的乘积的同构。(A, A)-双模 \(a_{i}\) 与 \(A/b_{i}\) 同构，故 \(A/b_{i}\) 的每个双边理想都等于 0 或 \(A/b_{i}\)。若环 A 是左阿廷的，则诸环 \(A/b_{i}\) 亦然，从而它们都是单的（VIII, 第 120 页，定义 1）。

最后证明 (iv) 蕴含 (i)。设 A 是一个有限族 \((\mathrm{A}_{i})_{i\in\mathrm{I}}\) 单环的乘积。用 \(\pi_{i}\) 表示从 A 到 \(A_{i}\) 的指标为 i 的投影，用 \(M_{i}\) 表示以 \(A_{i}\) 为底层加群、以 \((a,x)\mapsto\pi_{i}(a)x\) 为作用律的 A-模。由于环 \(A_{i}\) 是单的，\(A_{i}\)-模 \((\mathrm{A}_{i})_{s}\) 是半单的，故 A-模 \(M_{i}\) 是半单的。由于 A-模 \(A_{s}\) 无非就是乘积 \(\prod_{i\in I}M_{i}\)，它是半单的。

定义 1. —— 若环 A 具有定理 1 的等价性质 (i) 至 (vi)，则称环 A 是半单的。交换环 k 上的代数 A，若其底层环是半单的，则称为半单代数。

命题 1. —— 设 A 是一个半单环。存在 A 的极小左理想的一个有限族 \((\mathfrak{m}_{i})_{i\in I}\)，使得 \(A_{s}=\oplus_{i\in I}m_{i}\)。若 \((\mathfrak{m}_{i})_{i\in I}\) 是这样的族，则每个单 A-模都同构于某个 \(m_{i}\)。单 A-模的类所成的集是有限的。

第一个断言由 A-模 \(A_{s}\) 是半单的且有限生成这一事实得出。每个单模都同构于 \(A_{s}\) 的一个商（VIII, 第 46 页，命题 1）。第二个断言随即由 VIII, 第 56 页的推论 3 得出，而第三个断言是第二个的直接推论。

例. —— 设 G 是一个有限群，K 是一个交换域。*我们将进一步看到（VIII, 第 401 页，推论 1），体 K 上的群 G 的代数 K{[}G{]} 是半单环，当且仅当 K 的特征指数与 G 的阶互素。*

注. —— 1) 设 K 是一个交换域，A 是 K 上的一个半单代数。则存在为体的 K-代数 \(D_{1},\ldots,D_{s}\) 与整数 \(r_{1}\geqslant1,\ldots,r_{s}\geqslant1\)，使得 K-代数 A 同构于乘积 \(\prod_{i=1}^{s}\mathbf{M}_{r_{i}}(\mathrm{D}_{i})\)。

\begin{enumerate}
\def\labelenumi{\arabic{enumi})}
\setcounter{enumi}{1}
\tightlist
\item
  设 K 是一个代数闭域，A 是 K 上的一个有限次数代数。由 VIII, 第 122 页的注 4，代数 A 是半单的当且仅当存在整数 \(n_{1} \geqslant 1, \ldots, n_{r} \geqslant 1\)，使得 A 同构于代数 \(\mathbf{M}_{n_{1}}(\mathrm{K}) \times \cdots \times \mathbf{M}_{n_{r}}(\mathrm{K})\)。
\end{enumerate}

命题 2. —— a) 半单环的中心是半单的。

\begin{enumerate}
\def\labelenumi{\alph{enumi})}
\setcounter{enumi}{1}
\item
  半单环的反环是半单的。
\item
  半单环对双边理想的商是半单环。
\item
  有限族半单环的乘积是半单环。
\end{enumerate}

设 A 是一个半单环。它同构于一个有限族 \((\mathrm{A}_{i})_{i\in\mathrm{I}}\) 单环的乘积。A 的中心同构于诸 \(A_{i}\) 的中心的乘积，而 \(A^{o}\) 同构于诸 \(A_{i}^{o}\) 的乘积环。于是断言 a) 与 b) 由 VIII, 第 121 页的推论 1 得出。

设 \(\mathfrak{a}\) 是 A 的一个双边理想。A-模 \(A_s / \mathfrak{a}\) 是半单 A-模 \(A_s\) 的商，故是半单的。于是 \((A / \mathfrak{a})\)-模 \(A_s / \mathfrak{a}\) 是半单的；由此得断言 c)。

一个环是半单的当且仅当它同构于一个有限族单环的乘积；由此得断言 d)。

命题 3. —— 设 A 是一个交换环。下列性质等价：

\begin{enumerate}
\def\labelenumi{(\roman{enumi})}
\item
  环 A 是半单的。
\item
  环 A 是阿廷的且约化的（V, §6, No.~7, 第 34 页）。
\item
  环 A 同构于一个有限族交换域的乘积。
\end{enumerate}

交换单环是交换域（VIII, 第 120 页，注 1）。故 (i) 等价于 (iii)。

显然 (iii) 蕴含 (ii)。反之，设环 A 是阿廷的且约化的。A 的素理想所成之集的交由 A 的幂零元组成（V, §15, No.~1, 第 118 页，命题 2），由于 A 是约化的，它约化为 0。由 VIII, 第 2 页，存在 A 的互异的素理想 \(\mathfrak{p}_1, \ldots, \mathfrak{p}_r\)，使得我们有 \(\mathfrak{p}_1 \cap \cdots \cap \mathfrak{p}_r = 0\)。由 VIII, 第 8 页的推论，阿廷环 A 的每个素理想 \(\mathfrak{p}_i\) 都是极大的；于是我们有 \(\mathfrak{p}_i + \mathfrak{p}_j = A\)（只要 \(i\) 与 \(j\) 互异）。由 I, §8, No.~11, 第 110 页的命题 9，从 A 到环 \(\prod_{i=1}^{r} (A / \mathfrak{p}_i)\) 的典范同态是一个同构。环 \(A / \mathfrak{p}_i\) 对每个 \(i\) 都是域，故 (ii) 蕴含 (iii)。

体上的有限次数交换代数是阿廷交换环。于是命题 3 推广了 V, §6, No.~7, 第 34 页的命题 5。

